% GENERADO por docs/libros/arte/desde_app.js desde js/explora-*.js: no editar a mano.
\rutade{pu-cm-2}{{Partes iguales},{Fracciones equivalentes},{Comparar y sumar},{La fracción de una cantidad},{Ahora sí, un problema}}
\andamiode{pu-cm-2}{\item El denominador dice en cuántas partes \textbf{iguales} se corta; el numerador, cuántas se toman.
  \item Multiplicar o dividir arriba y abajo por el mismo número da una fracción \textbf{equivalente}: la misma cantidad.
  \item Para comparar o sumar, primero partes del mismo tamaño (el mismo denominador); luego se cuentan.
  \item La fracción de una cantidad: se divide entre el denominador y se multiplica por el numerador.}
\rutade{pu-cm-3}{{De cada 100},{Fracción, decimal y porcentaje},{El porcentaje de una cantidad},{Subir y bajar},{Ahora sí, un problema}}
\andamiode{pu-cm-3}{\item Por ciento es «de cada 100»: el $100\,\%$ es el total.
  \item Fracción, decimal y porcentaje son tres formas de lo mismo: $\tfrac15 = 0.2 = 20\,\%$.
  \item El porcentaje de una cantidad: el porcentaje en decimal por la cantidad. El $10\,\%$ corre el punto.
  \item Bajar un porcentaje es multiplicar por menos de 1; subir, por más de 1. Los cambios seguidos se multiplican.}
\rutade{pu-cm-4}{{Multiplicar repetido},{Al cuadrado: un área},{Contar factores},{Raíces},{Ahora sí, un problema}}
\andamiode{pu-cm-4}{\item Una potencia es multiplicar repetido: $b^n$ son $n$ factores iguales a $b$. Y $b^0 = 1$.
  \item Elevar al cuadrado es el área de un cuadrado: $n^2 = n \cdot n$.
  \item Con la misma base, al multiplicar potencias los exponentes se suman; al elevar una potencia, se multiplican.
  \item La raíz cuadrada busca el lado: $\sqrt{81} = 9$ porque $9^2 = 81$. Si no es exacta, se acorrala entre dos cuadrados.}
\rutade{pu-cm-1}{{La recta},{Sumar y restar},{El menos da la vuelta},{Las capas},{Ahora sí, un problema}}
\andamiode{pu-cm-1}{\item A la izquierda del 0 están los negativos; a la derecha, los positivos. Es menor el que está más a la izquierda.
  \item Un número y su opuesto están igual de lejos del 0 y suman 0. Dos menos se cancelan: $-(-6) = 6$.
  \item La distancia es cuántos pasos separan dos números. Se escribe con barras y nunca es negativa.
  \item Sumar es moverse: un positivo lleva a la derecha y un negativo, a la izquierda.
  \item Restar es sumar el opuesto. Restar un negativo sube: $3 - (-4) = 3 + 4 = 7$.
  \item Multiplicar es repetir saltos. Multiplicar por un negativo \textbf{da la vuelta} en el 0: $(-3)(-4) = 12$.
  \item Número \textbf{par} de signos menos: positivo. Número \textbf{impar}: negativo.
  \item Primero paréntesis, luego potencias, después \ensuremath{\times} y \ensuremath{\div}, y al final + y \ensuremath{-}.}
\rutade{pu-ge-1}{{Ángulos que se completan},{Paralelas y transversal},{Los ángulos de un triángulo},{Polígonos: cortar en triángulos},{Ahora sí, un problema}}
\andamiode{pu-ge-1}{\item Vecinos sobre una recta suman 180\textdegree{}; en una esquina recta, 90\textdegree{}; los opuestos por el vértice son iguales.
  \item Con paralelas: correspondientes y alternos internos son iguales; conjugados internos suman 180\textdegree{}.
  \item Los tres ángulos de un triángulo suman 180\textdegree{}, y un exterior es la suma de los dos interiores lejanos.
  \item Un polígono de n lados se corta en n \ensuremath{-} 2 triángulos: sus ángulos suman $(n - 2) \cdot 180^\circ$.}
\rutade{pu-ge-2}{{Tres cuadrados},{El lado que falta},{Triángulos que conviene reconocer},{Ahora sí, un problema}}
\andamiode{pu-ge-2}{\item En un triángulo rectángulo, el cuadrado de la hipotenusa vale lo que los dos de los catetos: $c^2 = a^2 + b^2$.
  \item Hipotenusa: se suman los cuadrados. Cateto: se restan. Al final, la raíz.
  \item Reconoce 3-4-5 y sus múltiplos, el medio cuadrado $x, x, x\sqrt2$ y el medio equilátero $x, x\sqrt3, 2x$.}
\rutade{pu-ge-3}{{La misma forma},{Sombras y Tales},{Largo, área y volumen},{Ahora sí, un problema}}
\andamiode{pu-ge-3}{\item Dos figuras son semejantes si tienen los mismos ángulos y todos los lados se multiplican por el \textbf{mismo} número, la razón.
  \item Con los rayos paralelos, altura y sombra guardan la misma razón en todos los objetos.
  \item Si los largos se multiplican por k, las áreas por $k^2$ y los volúmenes por $k^3$.}
\rutade{pu-ge-4}{{Base por altura},{El rombo: media caja},{El trapecio: dos juntos},{El círculo},{El sector: una tajada},{Por el borde},{Ahora sí, un problema}}
\andamiode{pu-ge-4}{\item Rectángulo y paralelogramo: $b \cdot h$. Triángulo: la mitad. La altura siempre es perpendicular a la base.
  \item Rombo: la mitad del producto de las diagonales, $\tfrac{D \cdot d}{2}$.
  \item Área del trapecio: promedio de las bases por la altura, $\tfrac{(B + b)\,h}{2}$.
  \item El borde mide $2\pi r$ y el círculo ocupa $\pi r^2$. Cuidado con usar el diámetro donde va el radio.
  \item Sector y arco: la fracción $\tfrac{\theta}{360^\circ}$ del área y del borde del círculo.
  \item Para el borde de una figura armada, recórrela por fuera: los lados que comparten las piezas no cuentan.}
\rutade{pu-ge-5}{{La grande menos lo que falta},{Triángulos en la cuadrícula},{Mover piezas},{Ahora sí, un problema}}
\andamiode{pu-ge-5}{\item Área sombreada = figura grande \ensuremath{-} lo que se le quita. No hay que inventar una fórmula nueva.
  \item Figura torcida: enciérrala en un rectángulo y réstale los triángulos de las esquinas.
  \item Antes de calcular, busca si una pieza que sobra llena un hueco igual: el área no cambia al moverla.}
\rutade{pu-ge-6}{{Capas que se apilan},{Conos y pirámides: un tercio},{La piel: desarrollar la caja},{Ahora sí, un problema}}
\andamiode{pu-ge-6}{\item Prismas y cilindros: volumen = área de la base \ensuremath{\times} altura, porque son capas iguales.
  \item Cono y pirámide: un tercio de base por altura, porque tres llenan el cilindro o el prisma.
  \item Superficie = suma de las caras. En una caja van por parejas; en un cubo son $6a^2$.}
\rutade{pu-or-1}{{Dibuja la fila},{Bloques que viajan juntos},{Probar y descartar},{Ahora sí, un problema}}
\andamiode{pu-or-1}{\item Dibuja casillas numeradas y pon primero las pistas que dicen un puesto exacto.
  \item Junta en un bloque lo que va pegado y busca dónde cabe entero.
  \item Con dos posibilidades, supón una; si choca con una pista, la otra es la verdadera.}
\rutade{pu-or-2}{{Fija a uno},{Frente a frente},{Dar la vuelta contando},{Ahora sí, un problema}}
\andamiode{pu-or-2}{\item En círculo, fija a una persona: girar la mesa no crea un acomodo nuevo.
  \item Con un número par de sillas, «frente a» es avanzar la mitad de ellas.
  \item Contando en ronda: suma los pasos y, si te pasas, resta el número de sillas.}
\rutade{pu-or-3}{{Tachar y marcar},{Un \ding{51} que tacha a los demás},{Dos tablas que se hablan},{Ahora sí, un problema}}
\andamiode{pu-or-3}{\item Cada «no» es una \ding{55}. Cuando en una fila o columna queda un solo hueco, ese es el \ding{51}.
  \item Empieza por las pistas que afirman: su \ding{51} tacha fila y columna y hace aparecer los demás.
  \item Una pista que une dos datos («el piloto vive en Tunja») pasa las marcas de una tabla a la otra.}
\rutade{pu-or-4}{{Leer bien un «si…»},{«Solo si» y la vuelta},{Tiene que, podría, no puede},{Ahora sí, un problema}}
\andamiode{pu-or-4}{\item «Si A, entonces B» solo prohíbe un caso: que esté A sin B.
  \item «C solo si D» es «si C, entonces D». Y dado vuelta: sin D no hay C.
  \item Tiene que: en todos los casos. Podría: en alguno. No puede: en ninguno.}
\rutade{pu-pr-1}{{Tantos por cada tantos},{Repartir por partes},{Escalas},{Ahora sí, un problema}}
\andamiode{pu-pr-1}{\item Dos razones iguales forman una proporción: se multiplica (o divide) por lo mismo, y los productos en cruz coinciden.
  \item Suma las partes, divide el total entre ellas y multiplica por las de cada uno.
  \item Escala 1 : N: lo real es N veces lo del mapa, en la misma unidad. Luego se cambia de unidad.}
\rutade{pu-pr-2}{{Directa: crecen a la vez},{Inversa: una sube, la otra baja},{Compuesta: cuenta el trabajo},{Ahora sí, un problema}}
\andamiode{pu-pr-2}{\item Directa: el cociente no cambia. Halla cuánto vale \textbf{uno} y multiplica.
  \item Inversa: el producto no cambia. Si una se duplica, la otra queda a la mitad.
  \item Compuesta: reduce a \textbf{uno solo en una hora} y luego multiplica por los cuadros que haya.}
\rutade{pu-pr-3}{{Distancia, velocidad y tiempo},{Velocidad media: no es el promedio},{Trabajar juntos: se suman las rapideces},{Ahora sí, un problema}}
\andamiode{pu-pr-3}{\item $d = v \cdot t$. Si dos se acercan, sus velocidades se suman; si uno persigue a otro, se restan.
  \item Velocidad media = distancia total \ensuremath{\div} tiempo total. No se promedian las velocidades.
  \item Suma lo que cada uno hace por hora. Juntos siempre tardan menos que el más rápido.}
\rutade{pu-pr-4}{{El tanto por ciento de algo},{¿Qué porcentaje es?},{Al revés: hallar el total},{Ahora sí, un problema}}
\andamiode{pu-pr-4}{\item El p \% de N es $\tfrac{p}{100} \cdot N$. A mano, ármalo con piezas del 10 \%, del 5 \% y del 1 \%.
  \item Qué porcentaje es A de B: $\tfrac{A}{B} \cdot 100$. El total siempre va abajo.
  \item Para volver al total: divide hasta el 10 \% o el 1 \% y sube al 100 \%.}
\rutade{pu-pr-5}{{Subir y bajar es multiplicar},{Uno tras otro},{Devolverse: IVA y precio de antes},{Ahora sí, un problema}}
\andamiode{pu-pr-5}{\item Aumentar r \% es multiplicar por $1 + \tfrac{r}{100}$; descontar, por $1 - \tfrac{r}{100}$.
  \item Porcentajes sucesivos: multiplica los factores. Cada cambio se calcula sobre el valor nuevo.
  \item Adelante se multiplica por el factor; atrás se divide. Nunca se quita el mismo porcentaje.}
\rutade{pu-co-1}{{El árbol de opciones},{¿Se suma o se multiplica?},{Casilla por casilla},{Ahora sí, un problema}}
\andamiode{pu-co-1}{\item Si una decisión tiene m opciones y la otra n, juntas tienen $m \cdot n$.
  \item Una cosa \textbf{o} la otra: se suma. Una cosa \textbf{y} después la otra: se multiplica.
  \item Una casilla por decisión, primero la más exigente, y se multiplican.}
\rutade{pu-co-2}{{Ordenar a todos},{Solo algunos puestos},{Letras repetidas},{Los que van juntos},{Ahora sí, un problema}}
\andamiode{pu-co-2}{\item Ordenar n cosas distintas: $n! = n \cdot (n-1) \cdots 2 \cdot 1$ maneras.
  \item Con k puestos se multiplican k números que bajan de uno en uno desde n.
  \item Con repetidas: el factorial del total entre el factorial de cada grupo de iguales.
  \item Ata a los que van juntos, ordena los paquetes y multiplica por los órdenes de adentro.}
\rutade{pu-co-3}{{El orden no importa},{Dividir entre los órdenes},{Elegir es también dejar},{Ahora sí, un problema}}
\andamiode{pu-co-3}{\item Si el orden no importa, es una combinación: se cuentan los órdenes y se quitan los repetidos.
  \item Combinaciones de n en k: $\binom{n}{k} = \tfrac{n (n-1) \cdots}{k!}$, con k factores arriba.
  \item $\binom{n}{k} = \binom{n}{n-k}$: cuenta el grupo más pequeño, el que va o el que se queda.}
\rutade{pu-co-4}{{Sumar sin contar doble},{Contar lo contrario},{Separar en casos},{Ahora sí, un problema}}
\andamiode{pu-co-4}{\item Uno u otro: suma los dos grupos y resta los que están en ambos.
  \item Al menos uno = todos \ensuremath{-} ninguno.
  \item Divide en casos que no se crucen, cuenta cada uno y suma.}
\rutade{pu-co-5}{{Contar por tamaños},{Diagonales y parejas},{Caminos que se suman},{Ahora sí, un problema}}
\andamiode{pu-co-5}{\item Cuenta por tamaños: para cada uno, en cuántos lugares cabe. Luego suma.
  \item Un polígono de n lados tiene $\tfrac{n (n - 3)}{2}$ diagonales.
  \item Los caminos a una esquina son los de su izquierda más los de abajo.}
\rutade{pu-es-1}{{Caminar por las calles},{Juntar los movimientos},{En línea recta},{Ahora sí, un problema}}
\andamiode{pu-es-1}{\item Distancia por calles = cuadras a lo ancho + cuadras a lo alto.
  \item Norte contra sur, este contra oeste: se restan. Lo que queda dice dónde terminas.
  \item Por calles se suman los lados; en línea recta se usa $\sqrt{a^2 + b^2}$.}
\rutade{pu-es-2}{{Contar por torres},{Las vistas},{El cubo pintado},{Ahora sí, un problema}}
\andamiode{pu-es-2}{\item Escribe la altura de cada torre y suma: así cuentas también los cubos escondidos.
  \item En una vista, cada columna muestra la altura \textbf{máxima} de las torres que están en esa dirección.
  \item Con arista n: 8 esquinas, $12(n-2)$ en aristas, $6(n-2)^2$ en caras y $(n-2)^3$ sin pintar.}
\rutade{pu-es-3}{{Caras opuestas en el molde},{El dado suma siete},{Rodar el dado},{Ahora sí, un problema}}
\andamiode{pu-es-3}{\item En una fila recta, la cara opuesta está \textbf{saltando una}. Las dos que sobran son opuestas.
  \item Opuestas: 1 con 6, 2 con 5 y 3 con 4. Las seis caras suman 21.
  \item Al volcar a la derecha: arriba \ensuremath{\to} derecha \ensuremath{\to} abajo \ensuremath{\to} izquierda \ensuremath{\to} arriba. El frente y atrás no cambian.}
\rutade{pu-es-4}{{El espejo},{Los giros},{Doblar y perforar},{Ahora sí, un problema}}
\andamiode{pu-es-4}{\item Reflejar: mismo renglón y misma distancia al espejo, del otro lado.
  \item Girar da vueltas; reflejar voltea. 270\textdegree{} a un lado es 90\textdegree{} al otro.
  \item Con d dobleces hay $2^d$ capas: cada perforación deja $2^d$ huecos.}
\rutade{pu-cd-1}{{El alfabeto corrido},{Letras con valor},{Bombillos: el binario},{Ahora sí, un problema}}
\andamiode{pu-cd-1}{\item Cifrar: avanzar k lugares. Descifrar: devolverse k. Después de Z sigue A.
  \item Memoriza E 5, J 10, O 15, T 20 y cuenta desde la más cercana.
  \item Empieza por el bombillo más grande que quepa y sigue con lo que falta.}
\rutade{pu-cd-2}{{La máquina de números},{Letras que cambian de puesto},{El corrimiento que crece},{Ahora sí, un problema}}
\andamiode{pu-cd-2}{\item Propón una regla con un ejemplo y pruébala con \textbf{todos}. Si falla en uno, se descarta.
  \item Si las letras son las mismas, numera los puestos y mira a dónde va cada uno.
  \item Escribe cuánto avanza cada letra: si no es igual, busca el patrón de los avances.}
\rutade{pu-gr-1}{{Leer el eje},{¿Cuánto cambió, en porcentaje?},{La torta},{Ahora sí, un problema}}
\andamiode{pu-gr-1}{\item Primero la escala del eje, luego la punta de cada barra.
  \item Variación porcentual = cambio \ensuremath{\div} valor de \textbf{antes}.
  \item El 100 \% son 360\textdegree{}, así que cada 1 \% son 3.6\textdegree{}.}
\rutade{pu-gr-2}{{Promedio: nivelar},{Mediana y moda},{No todos pesan igual},{Ahora sí, un problema}}
\andamiode{pu-gr-2}{\item Promedio = suma de todos \ensuremath{\div} cuántos son. Es el nivel donde todo queda parejo.
  \item Mediana: ordena y toma el del medio (con cantidad par, el promedio de los dos del medio). Moda: el más repetido.
  \item Promedio de grupos: (valor \ensuremath{\times} cuántos) de cada grupo, sumado, entre el total de personas.}
\rutade{pu-gr-3}{{El eje que engaña},{Cantidad no es proporción},{¿A quién se le preguntó?},{Ahora sí, un problema}}
\andamiode{pu-gr-3}{\item Antes de comparar alturas, mira dónde empieza el eje. Compara los números, no el dibujo.
  \item La cantidad cuenta cuántos; la proporción dice cuántos \textbf{de cada cien}.
  \item Una muestra representa a todos si se escoge al azar y de todas partes.}
\rutade{pu-se-1}{{Siempre se suma lo mismo},{Los saltos también cambian},{Otras reglas},{Ahora sí, un problema}}
\andamiode{pu-se-1}{\item Término n = primero + (n \ensuremath{-} 1) saltos. Los saltos son uno menos que los términos.
  \item Escribe las diferencias. Si no son iguales, busca la regla de las diferencias.
  \item Prueba: ¿se suma?, ¿se multiplica?, ¿se suman los dos anteriores?, ¿son dos mezcladas?}
\rutade{pu-se-2}{{Las letras son números},{Figuras que crecen},{Dos reglas a la vez},{Ahora sí, un problema}}
\andamiode{pu-se-2}{\item Cambia las letras por su lugar en el alfabeto y busca la regla de los números.
  \item Busca la parte fija y la que crece: aquí 1 fijo y 4 más en cada figura, $4n + 1$.
  \item Separa las letras de los números, resuelve cada regla y junta las respuestas.}
\rutade{pu-se-3}{{Lo que se repite, sobra},{La última cifra},{Dos ciclos a la vez},{Ahora sí, un problema}}
\andamiode{pu-se-3}{\item Divide entre el largo del ciclo: el \textbf{residuo} dice el lugar. Residuo 0 es el último.
  \item Escribe las últimas cifras hasta que se repitan y usa el residuo del exponente.
  \item Dos ciclos vuelven a coincidir en el mínimo común múltiplo de sus largos.}
\rutade{pu-lo-1}{{Y, o, no: tres circuitos},{La tabla de verdad},{El condicional: una promesa},{Ahora sí, un problema}}
\andamiode{pu-lo-1}{\item «y»: verdadera solo con ambas. «o»: falsa solo con ninguna. «no»: lo contrario.
  \item Una fila por caso: con n letras, $2^n$ filas. Se llena por partes.
  \item $p \to q$ solo es falsa cuando p es verdadera y q es falsa.}
\rutade{pu-lo-2}{{Negar es decir lo contrario, completo},{Comprobarlo con la tabla},{Todos, alguno, ninguno},{Ahora sí, un problema}}
\andamiode{pu-lo-2}{\item De Morgan: niega cada parte y cambia «y» por «o» o «o» por «y».
  \item Si las columnas coinciden en todas las filas, las frases son equivalentes.
  \item La negación de «todos son» es «alguno no es». La de «alguno es», «ninguno es».}
\rutade{pu-lo-3}{{Uno dentro del otro},{La vuelta que sí vale},{Las dos trampas},{Ahora sí, un problema}}
\andamiode{pu-lo-3}{\item «Si p, entonces q»: p está dentro de q. p basta; q hace falta.
  \item $p \to q$ equivale a $\neg q \to \neg p$: negar las dos y voltear.
  \item Seguras: de p se llega a q, y de «no q» a «no p». Con q o con «no p» no se concluye nada.}
\rutade{pu-lo-4}{{Todos y ninguno},{Algunos},{Encadenar frases},{Ahora sí, un problema}}
\andamiode{pu-lo-4}{\item «Todos los A son B»: A dentro de B. «Ningún A es B»: separados.
  \item «Algunos A son B»: hay al menos uno en la parte común. Se puede voltear.
  \item Dibuja las premisas. Una conclusión vale si sale en todos los dibujos posibles.}
\rutade{pu-lo-5}{{Suponer y mirar si choca},{¿Quién fue? Una tabla de casos},{Atajos},{Ahora sí, un problema}}
\andamiode{pu-lo-5}{\item Supón una cosa; si choca, la contraria es la verdadera. Al final, comprueba todo.
  \item Una fila por sospechoso: escribe V o F en cada frase y busca la fila que cumple la condición.
  \item Frases opuestas: una es verdad y la otra mentira, sin saber cuál. Eso decide a los demás.}
\rutade{pu-op-1}{{Una receta con huecos},{Los negativos, con paréntesis},{Operaciones en una tabla},{Ahora sí, un problema}}
\andamiode{pu-op-1}{\item Reemplaza cada número en el lugar de su letra. El primero es a; el segundo, b.
  \item Los negativos entran con paréntesis: $(-3)^2 = 9$ y restar $-4$ es sumar 4.
  \item Fila del primero, columna del segundo. El neutro no cambia a nadie.}
\rutade{pu-op-2}{{De adentro hacia afuera},{Cuando hay una x},{Dos operaciones juntas},{Ahora sí, un problema}}
\andamiode{pu-op-2}{\item Resuelve el paréntesis, ponlo como un solo número y sigue.
  \item Reemplaza, y resuelve la ecuación que queda. Al final, comprueba.
  \item Cada símbolo, su receta; primero el de adentro.}
\rutade{pu-ar-1}{{El segundo pedazo es de lo que queda},{Multiplicar lo que queda},{Al revés: hallar el principio},{Ahora sí, un problema}}
\andamiode{pu-ar-1}{\item Después de cada gasto hay un todo nuevo: la fracción siguiente se toma de ahí.
  \item Si se gasta $\tfrac16$ queda $\tfrac56$; si luego se gasta $\tfrac14$ de eso, queda $\tfrac34 \cdot \tfrac56$.
  \item De atrás hacia adelante: lo que sobró es una fracción conocida de lo que había antes.}
\rutade{pu-ar-2}{{La suma escondida},{Entra uno, sale uno},{Notas con peso},{Ahora sí, un problema}}
\andamiode{pu-ar-2}{\item Suma = cuántos son \ensuremath{\times} promedio. Con la suma se resuelven los problemas de promedios.
  \item Suma vieja \ensuremath{\pm} el dato = suma nueva; divide entre la cantidad nueva.
  \item Promedio ponderado: suma de nota \ensuremath{\times} peso, con los pesos en decimales.}
\rutade{pu-ar-3}{{Cada cifra según su puesto},{Columna por columna},{Letras repetidas},{Ahora sí, un problema}}
\andamiode{pu-ar-3}{\item $\overline{AB} = 10A + B$. Al voltear las cifras, cambia en 9 veces su diferencia.
  \item Primero unidades; si pasan de 9, se lleva 1. Luego decenas, sumando lo que se llevó.
  \item Escribe cada número con 10 y 1, junta las letras iguales y saca el factor común.}
\rutade{pu-ar-4}{{El peor caso},{El producto más grande},{Repartos extremos},{Ahora sí, un problema}}
\andamiode{pu-ar-4}{\item Imagina la peor suerte posible y suma 1.
  \item Si la suma es fija, el producto es mayor cuando los dos números son iguales.
  \item Máximo de uno: los demás al mínimo. Mínimo del mayor: todos lo más parejo posible.}
\rutade{pu-ec-1}{{Traducir palabra por palabra},{La balanza},{Todo con la misma x},{Ahora sí, un problema}}
\andamiode{pu-ec-1}{\item Ponle x a lo que no sabes; cada palabra es una operación, y «es» es el igual.
  \item Quita lo mismo a los dos lados hasta dejar sola la x; luego divide.
  \item Llama x a uno, escribe los demás con esa x y súmalos.}
\rutade{pu-ec-2}{{Las edades crecen juntas},{Cambiar una por otra},{Mezclas},{Ahora sí, un problema}}
\andamiode{pu-ec-2}{\item A los dos se les suman los mismos años. La diferencia es siempre la misma.
  \item Supón que todas valen lo menos; cada cambio sube la diferencia entre las dos.
  \item Valor de la mezcla = (cantidad \ensuremath{\times} precio de cada parte, sumado) \ensuremath{\div} cantidad total.}
\rutade{pu-ec-3}{{Dos recibos},{Suma y diferencia},{Desigualdades},{Ahora sí, un problema}}
\andamiode{pu-ec-3}{\item Si dos ecuaciones comparten una parte, réstalas: esa parte desaparece.
  \item Menor = (suma \ensuremath{-} diferencia) \ensuremath{\div} 2. Mayor = (suma + diferencia) \ensuremath{\div} 2.
  \item Se despeja como una ecuación, pero al multiplicar o dividir por un negativo se voltea el signo.}
\rutade{pu-cj-1}{{Cada cosa en su zona},{Las operaciones se pintan},{Contar sin repetir},{Ahora sí, un problema}}
\andamiode{pu-cj-1}{\item Pregunta dos veces: ¿está en A?, ¿está en B? Las respuestas dicen la zona.
  \item Unión: todo. Intersección: lo común. Diferencia: quitar lo común. Complemento: lo que falta del universo.
  \item $n(A \cup B) = n(A) + n(B) - n(A \cap B)$.}
\rutade{pu-cj-2}{{Del centro hacia afuera},{Cuando no dicen el centro},{Tres conjuntos},{Ahora sí, un problema}}
\andamiode{pu-cj-2}{\item Primero los que hacen las dos; luego se restan; al final, los de afuera.
  \item Los dos = (A + B) \ensuremath{-} los que tienen alguno.
  \item Centro, luego las parejas (restando el centro) y al final los que hacen una sola.}
\rutade{pu-nu-1}{{Primos: los que no se acomodan},{Descomponer en primos},{Criterios de divisibilidad},{Ahora sí, un problema}}
\andamiode{pu-nu-1}{\item Primo: solo se acomoda en una fila. Compuesto: forma un rectángulo.
  \item Se divide entre primos hasta llegar a 1. Divisores: suma 1 a cada exponente y multiplica.
  \item 3 y 9: suma de cifras. 2, 5 y 10: última cifra. 4: dos últimas. 6: entre 2 y entre 3.}
\rutade{pu-nu-2}{{MCM: cuándo coinciden},{MCD: el pedazo más grande},{Con factores primos},{Ahora sí, un problema}}
\andamiode{pu-nu-2}{\item MCM: el menor número que es múltiplo de los dos. Sirve para «¿cuándo vuelven a coincidir?».
  \item MCD: el mayor número que divide a los dos. Sirve para repartir o cortar sin que sobre.
  \item MCD: primos comunes. MCM: todos los primos, con el mayor exponente.}
\rutade{pu-nu-3}{{Par o impar},{Los residuos se suman y se multiplican},{La última cifra},{Ahora sí, un problema}}
\andamiode{pu-nu-3}{\item Suma: impar con impar da par. Producto: solo impar por impar da impar.
  \item Cambia cada número por su residuo, haz la cuenta con ellos y saca el residuo otra vez.
  \item Para la última cifra de sumas, productos y potencias basta usar las últimas cifras.}
\rutade{pu-pb-1}{{Favorables sobre posibles},{Dos dados: la tabla},{De imposible a seguro},{Ahora sí, un problema}}
\andamiode{pu-pb-1}{\item $P = \tfrac{\text{casos favorables}}{\text{casos posibles}}$, con casos igual de fáciles.
  \item Dos dados: 36 parejas ordenadas. No todas las sumas son igual de probables.
  \item 0 es imposible, 1 es seguro, y ninguna probabilidad pasa de 1.}
\rutade{pu-pb-2}{{El complemento},{Independientes: se multiplican},{Sin devolver: cambia la bolsa},{Ahora sí, un problema}}
\andamiode{pu-pb-2}{\item $P(\text{no A}) = 1 - P(A)$. «Al menos uno» es 1 menos «ninguno».
  \item Si A y B son independientes, $P(A \text{ y } B) = P(A) \cdot P(B)$.
  \item Sin devolver: en el segundo paso hay una bola menos, abajo y (si salió de ese color) arriba.}
\rutade{pu-el-1}{{Los bordes del texto},{La fuente dice para qué},{La vista rápida},{Ahora sí, un problema}}
\andamiode{pu-el-1}{\item El título anuncia el tema; el autor, quién habla; la fuente, para qué; la fecha, de cuándo es.
  \item Un informe registra, la divulgación explica, una columna defiende y un testimonio cuenta.
  \item El comienzo de cada párrafo y las palabras que se repiten dicen dónde está cada cosa.}
\rutade{pu-el-2}{{Una idea por párrafo},{El mapa del texto},{Lo que sostiene y lo que adorna},{Ahora sí, un problema}}
\andamiode{pu-el-2}{\item La idea reúne lo que dicen las demás: ni un detalle suelto ni más de lo que dice el párrafo.
  \item «Sin embargo» opone, «además» suma, «por eso» concluye: el comienzo dice qué hace el párrafo.
  \item La prueba: quita la oración. Si la idea central pierde apoyo, era importante; si no, era adorno.}
\rutade{pu-el-3}{{El tiempo},{Volver al párrafo},{La intrusa disfrazada},{Ahora sí, un problema}}
\andamiode{pu-el-3}{\item Leer una vez, responder volviendo al texto, marcar la difícil y no dejar nada en blanco.
  \item En EXCEPTO cada opción busca su párrafo. La que no lo encuentra es la respuesta.
  \item Lo que está en otro párrafo no cuenta. Y deducir es sacar algo de lo escrito, sin inventar.}
\rutade{pu-tt-1}{{Cuatro maneras},{Las pistas},{Se mezclan},{Ahora sí, un problema}}
\andamiode{pu-tt-1}{\item Narrar es contar; describir, mostrar; explicar, decir cómo funciona; argumentar, defender.
  \item Pasos: explica. Acciones en pasado: narra. Sensaciones: describe. «Debería» o «porque» con un juicio: argumenta.
  \item Predomina la tipología que organiza todo el texto, no la que ocupa más renglones.}
\rutade{pu-tt-2}{{Quién habla y para qué},{Dos voces y sus signos},{A qué apunta una palabra},{Ahora sí, un problema}}
\andamiode{pu-tt-2}{\item Testimonio: lo vivido. Opinión: un juicio. Divulgación: ciencia. Crónica: hora y lugar. Entrevista: dos voces.
  \item La raya da el turno, los dos puntos anuncian y los de pregunta encierran lo que se pregunta.
  \item Para saber a qué apunta una palabra, pon cada opción en su lugar y lee toda la oración.}
\rutade{pu-tt-3}{{Cómo está armado},{La voz del testimonio},{Género o cita},{Ahora sí, un problema}}
\andamiode{pu-tt-3}{\item Argumentar: tesis, argumento, contraargumento y respuesta. Contar una vida: antes, cambio y después.
  \item Primera persona, marcas de oralidad como «uno» y palabras de tiempo que ordenan la vida.
  \item El género lo da quien habla en todo el texto. Una pregunta retórica no pide un dato: señala algo.}
\rutade{pu-ie-1}{{La palabra clave},{Datos que se juntan},{Ni más ni menos que lo escrito},{Ahora sí, un problema}}
\andamiode{pu-ie-1}{\item Saca una palabra clave de la pregunta (un nombre, una cifra) y vuelve al texto a buscarla.
  \item Si no está en una oración, mira si sale de juntar dos. Lo demás se deja quieto.
  \item «Pidió» no es «logró». Compara el verbo de la opción con el del texto.}
\rutade{pu-ie-2}{{Lo mismo, con otras palabras},{El alcance},{Las bisagras también dicen},{Ahora sí, un problema}}
\andamiode{pu-ie-2}{\item Una paráfrasis fiel no agrega ni quita: solo cambia las palabras.
  \item Cuida las palabras que miden: cuántos, qué tan seguido y hasta dónde.
  \item «Por eso» resulta, «sin embargo» limita y «además» suma: la paráfrasis conserva esa relación.}
\rutade{pu-ie-3}{{Verdadero no es «dicho»},{Los huecos que uno llena},{En otro lugar o en ninguno},{Ahora sí, un problema}}
\andamiode{pu-ie-3}{\item Para aceptar una opción, señala la frase que la dice. Si no puedes, no está.
  \item Que alguien lo haga, lo crea o lo pida no es que el texto lo afirme.
  \item Está aquí, está en otro lado o no está: se decide buscando, no recordando.}
\rutade{pu-ie-4}{{Cada idea tiene dueño},{La trampa de la voz},{Quién va antes y quién después},{Ahora sí, un problema}}
\andamiode{pu-ie-4}{\item «Según», las comillas y verbos como «sostiene» dicen de quién es la idea; «a mi juicio», que es del autor.
  \item No basta con que la idea esté: tiene que ser de quien dice la opción.
  \item Cuenta desde la fecha escrita; quien se menciona como conocido ya apareció; la causa va antes.}
\rutade{pu-in-1}{{Leer entre líneas},{Inferencia o suposición},{Gestos y cadenas},{Ahora sí, un problema}}
\andamiode{pu-in-1}{\item Una inferencia se apoya en pistas del texto: se puede señalar con el dedo.
  \item Si para sostenerla tienes que agregar algo tuyo, ya no es inferencia.
  \item «Desde entonces» encadena hechos; un gesto deja ver lo que se siente.}
\rutade{pu-in-2}{{Lo que se da por hecho},{Las palabras delatoras},{Dar por hecho no es decir},{Ahora sí, un problema}}
\andamiode{pu-in-2}{\item Presupuesto: lo que tiene que ser cierto para que la frase tenga sentido.
  \item Busca «por fin», «todavía», «ya no», «dejar de»… y también comparaciones, un «sí» que insiste o «saber».
  \item Lo que dice está escrito; lo que supone, tiene que ser cierto; lo demás no está.}
\rutade{pu-in-3}{{Encadenar hechos},{Después no es por},{Nombrar no es afirmar},{Ahora sí, un problema}}
\andamiode{pu-in-3}{\item De cada hecho pregunta qué tuvo que pasar antes. «Por eso» trae el efecto; «porque», la causa.
  \item Quita el primer hecho: si lo segundo pasaría igual, solo coinciden.
  \item «Hay quien culpa… pero»: el autor nombra una causa para negarla.}
\rutade{pu-in-4}{{Mostrar, no nombrar},{Los motivos},{Seguir el hilo del relato},{Ahora sí, un problema}}
\andamiode{pu-in-4}{\item Nostalgia, orgullo, resignación, desdén, reserva: cada una deja su huella en lo que se hace.
  \item El motivo se prueba con una frase del texto, no con lo que uno habría sentido.
  \item «Ese», «lo» y «hoy» dicen qué ya se contó; lo que sigue ayuda a decidir.}
\rutade{pu-vo-1}{{La prueba del reemplazo},{Verbos concretos, ideas},{Ni al pie de la letra ni de más},{Ahora sí, un problema}}
\andamiode{pu-vo-1}{\item Reemplaza, lee la oración completa y quédate con la que dice lo mismo.
  \item Piensa qué hace el verbo con una cosa y busca lo parecido con una idea.
  \item La justa recoge lo que el texto cuenta: ni literal ni exagerada.}
\rutade{pu-vo-2}{{Una palabra, varios sentidos},{El sentido más común},{Lo que retoma y lo que aclara},{Ahora sí, un problema}}
\andamiode{pu-vo-2}{\item Mira las palabras vecinas: ellas eligen el sentido.
  \item El sentido más común suele ser la trampa: pruébalo en la oración.
  \item «Ese…» retoma, los dos puntos aclaran y la coma puede abrir algo que se quita.}
\rutade{pu-vo-3}{{No es la suma de las palabras},{El contexto la explica},{Exagerar para decir},{Ahora sí, un problema}}
\andamiode{pu-vo-3}{\item No sumes las palabras: mira qué está pasando.
  \item Busca antes o después la situación que la explica.
  \item Imagínalo tal cual: si es imposible, es una manera de decir.}
\rutade{pu-vo-4}{{Las piezas},{Piezas que engañan},{El texto también ayuda},{Ahora sí, un problema}}
\andamiode{pu-vo-4}{\item Parte la palabra: comienzo, raíz y final. Busca otras que tengan la misma pieza.
  \item Prueba la lectura por piezas en la oración; si no encaja, manda el contexto.
  \item «Es decir» y los dos puntos traducen; «podría» no asegura; «en cambio» compara.}
\rutade{icf-ma-datos-1}{{Primero, la escala},{Cambio y porcentaje},{La torta},{Ahora sí, un problema}}
\andamiode{icf-ma-datos-1}{\item Mira cuánto vale una raya y qué parte de ella alcanza la barra.
  \item Cambio porcentual = lo que cambió \ensuremath{\div} lo que había antes.
  \item El sector es la misma fracción de 360\textdegree{} que la parte es del total.}
\rutade{icf-ma-datos-2}{{Repartir parejo},{El del centro y el que más sale},{El dato que jala y el que pesa},{Ahora sí, un problema}}
\andamiode{icf-ma-datos-2}{\item Promedio: suma todo y reparte entre cuántos son.
  \item Mediana: ordena y busca el centro. Moda: el que más se repite.
  \item Un dato extremo mueve el promedio y casi nunca la mediana.}
\rutade{icf-ma-datos-3}{{El eje que engaña},{Grupos de distinto tamaño},{La muestra},{Ahora sí, un problema}}
\andamiode{icf-ma-datos-3}{\item Antes de comparar barras, mira dónde empieza el eje.
  \item Cada cantidad se divide entre el total de su propio grupo.
  \item Sirve si cualquiera de la población podía salir.}
\rutade{icf-ma-proporcion-1}{{Comparar dividiendo},{Partes iguales},{La escala},{Ahora sí, un problema}}
\andamiode{icf-ma-proporcion-1}{\item Una razón se conserva al multiplicar o dividir las dos cantidades.
  \item Suma la razón, divide el total y multiplica por cada número.
  \item Del mapa a la realidad se multiplica; de la realidad al mapa, se divide.}
\rutade{icf-ma-proporcion-2}{{De cada cien},{Del pedazo al total},{Los cambios multiplican},{Ahora sí, un problema}}
\andamiode{icf-ma-proporcion-2}{\item Halla el 10 \% y arma con él cualquier porcentaje.
  \item Parte \ensuremath{\div} porcentaje (en decimal) = total.
  \item Subir p \%: por 1 + p. Bajar p \%: por 1 \ensuremath{-} p. Para volver, divide.}
\rutade{icf-ma-proporcion-3}{{Crecen juntas},{Una sube y la otra baja},{Lo que hace uno en una hora},{Ahora sí, un problema}}
\andamiode{icf-ma-proporcion-3}{\item Directa: el cociente no cambia. Halla lo que vale uno.
  \item Inversa: el producto no cambia. Más obreros, menos días.
  \item Halla lo que hace una en una hora y multiplica por lo que tengas.}
\rutade{icf-ma-proporcion-4}{{Algo por cada uno},{Distancia, velocidad y tiempo},{Juntos o detrás},{Ahora sí, un problema}}
\andamiode{icf-ma-proporcion-4}{\item Para comparar, halla cuánto corresponde a una sola unidad.
  \item d = v \ensuremath{\cdot} t; los minutos se pasan a fracción de hora.
  \item De frente: distancia \ensuremath{\div} suma. Detrás: ventaja \ensuremath{\div} diferencia.}
\rutade{icf-ma-azar-1}{{Y multiplica},{O suma},{Lo que se cuenta dos veces},{Ahora sí, un problema}}
\andamiode{icf-ma-azar-1}{\item Una cosa y además otra: se multiplica.
  \item Dos decisiones se multiplican; una sola elección entre grupos se suma.
  \item Alguno = suma de los dos grupos \ensuremath{-} los que están en ambos.}
\rutade{icf-ma-azar-2}{{¿Importa el orden?},{Con orden},{Sin orden se divide},{Ahora sí, un problema}}
\andamiode{icf-ma-azar-2}{\item Cambia el orden: si sale otra cosa, importa.
  \item Con orden: n por (n \ensuremath{-} 1) por… tantas veces como puestos.
  \item Sin orden: cuenta con orden y divide entre los órdenes del grupo.}
\rutade{icf-ma-azar-3}{{Favorables entre posibles},{De 0 a 1},{Lo contrario y dos dados},{Ahora sí, un problema}}
\andamiode{icf-ma-azar-3}{\item Favorables entre posibles; y se compara la parte, no la cantidad.
  \item Nunca es negativa ni pasa de 1.
  \item Contrario: 1 menos. Dos dados: 36 parejas ordenadas.}
\rutade{icf-ma-azar-4}{{¿Sobre quiénes?},{Con o sin devolver},{Al menos uno},{Ahora sí, un problema}}
\andamiode{icf-ma-azar-4}{\item Primero di entre quiénes se escoge: ese es el total.
  \item Sin devolver hay una menos en el total, y quizá una menos que sirve.
  \item Alguno = 1 \ensuremath{-} ninguno.}
\rutade{icf-ma-medida-1}{{Borde y superficie},{Medio rectángulo},{El promedio de las bases},{Ahora sí, un problema}}
\andamiode{icf-ma-medida-1}{\item Perímetro: se rodea, en metros. Área: se cubre, en metros cuadrados.
  \item Triángulo: base por altura entre 2.
  \item Trapecio: promedio de bases por altura. Rombo: diagonales entre 2.}
\rutade{icf-ma-medida-2}{{El círculo},{Partes y anillos},{Sumar piezas o quitar sobras},{Ahora sí, un problema}}
\andamiode{icf-ma-medida-2}{\item Borde 2\ensuremath{\pi}r, superficie \ensuremath{\pi}r\ensuremath{{}^2}. Una vuelta avanza el borde.
  \item Sector: fracción del círculo. Anillo: afuera menos adentro.
  \item Suma las piezas o quita lo que sobra: da igual.}
\rutade{icf-ma-medida-3}{{Por capas},{La tercera parte},{Piel y litros},{Ahora sí, un problema}}
\andamiode{icf-ma-medida-3}{\item Prisma o cilindro: área de la base por altura.
  \item Cono y pirámide: un tercio. Esfera: dos tercios de su cilindro.
  \item Piel: suma de caras. 1 m\ensuremath{{}^3} = 1 000 L; 1 L = 1 000 cm\ensuremath{{}^3}.}
\rutade{icf-ma-lineal-1}{{Por cada uno},{Sube, baja o nada},{Con unidades},{Ahora sí, un problema}}
\andamiode{icf-ma-lineal-1}{\item Pendiente: lo que sube entre lo que avanza, igual en toda la recta.
  \item Positiva sube, negativa baja, cero no cambia.
  \item Lo que cambia por cada unidad de lo que avanza.}
\rutade{icf-ma-lineal-2}{{Fijo más por cada uno},{En el plano},{Devolverse},{Ahora sí, un problema}}
\andamiode{icf-ma-lineal-2}{\item Total = fijo + cambio \ensuremath{\cdot} cantidad.
  \item Corta el eje en lo fijo y se inclina según el cambio.
  \item Quita lo fijo, reparte y comprueba.}
\rutade{icf-ma-lineal-3}{{El cruce},{De cada lado},{Decidir},{Ahora sí, un problema}}
\andamiode{icf-ma-lineal-3}{\item Empate: lo que se llevan entre lo que se acercan.
  \item Antes del cruce, menor fijo; después, menor pendiente.
  \item No hay plan mejor: hay plan mejor para tu cantidad.}
\rutade{icf-ma-forma-1}{{Media vuelta},{Tres ángulos},{¿Cierra?},{Ahora sí, un problema}}
\andamiode{icf-ma-forma-1}{\item Vecinos: suman 180\textdegree{}. Opuestos: iguales.
  \item Interiores: 180\textdegree{}. Exterior: la suma de los dos lejanos.
  \item El tercer lado va entre la resta y la suma de los otros.}
\rutade{icf-ma-forma-2}{{Cuadrados sobre los lados},{Falta un cateto},{Familias y atajos},{Ahora sí, un problema}}
\andamiode{icf-ma-forma-2}{\item Hipotenusa al cuadrado = suma de los catetos al cuadrado.
  \item Hipotenusa: suma. Cateto: resta.
  \item Reconoce 3-4-5, 5-12-13, 8-15-17 y sus múltiplos.}
\rutade{icf-ma-forma-3}{{Misma forma},{Paralelas},{Al escalar},{Ahora sí, un problema}}
\andamiode{icf-ma-forma-3}{\item Halla por cuánto creció un lado y úsalo en los demás.
  \item Tales: la razón entre los trozos es la misma.
  \item Largo: k. Área: k\ensuremath{{}^2}. Volumen: k\ensuremath{{}^3}.}
\rutade{icf-ma-forma-4}{{Tres nombres},{Solo importa el ángulo},{Medir sin llegar},{Ahora sí, un problema}}
\andamiode{icf-ma-forma-4}{\item Seno op/hip, coseno ady/hip, tangente op/ady.
  \item Cada ángulo tiene su seno, su coseno y su tangente.
  \item Escoge la razón que une lo que tienes con lo que buscas.}
\rutade{icf-ma-ecuaciones-1}{{La balanza},{De la frase a la letra},{Comprobar},{Ahora sí, un problema}}
\andamiode{icf-ma-ecuaciones-1}{\item Lo que haces en un lado lo haces en el otro.
  \item Una sola letra, y todo escrito con ella.
  \item Responde lo que preguntan y comprueba.}
\rutade{icf-ma-ecuaciones-2}{{Lo semejante se junta},{Un producto es un rectángulo},{Sacar lo común},{Ahora sí, un problema}}
\andamiode{icf-ma-ecuaciones-2}{\item x con x, x\ensuremath{{}^2} con x\ensuremath{{}^2}, números con números.
  \item Cuatro pedazos; las dos franjas se juntan.
  \item Saca lo que divide a todos y déjalo afuera.}
\rutade{icf-ma-ecuaciones-3}{{Comparar recibos},{Suponer que todos son iguales},{Suma y diferencia},{Ahora sí, un problema}}
\andamiode{icf-ma-ecuaciones-3}{\item Resta los recibos: solo queda lo distinto.
  \item Supón, mira lo que falta y reparte entre lo que sube cada cambio.
  \item Quita la ventaja, reparte y devuélvela.}
\rutade{icf-ma-ecuaciones-4}{{Muchas respuestas},{El signo que se voltea},{Como máximo, al menos},{Ahora sí, un problema}}
\andamiode{icf-ma-ecuaciones-4}{\item Pinta hacia donde se cumple y decide si el borde entra.
  \item Solo un negativo que multiplica o divide lo voltea.
  \item Máximo \ensuremath{\le}, al menos \ensuremath{\ge}; con «de» el borde no entra.}
\rutade{icf-ma-argumenta-1}{{Si… entonces},{Basta o hace falta},{Paso a paso},{Ahora sí, un problema}}
\andamiode{icf-ma-argumenta-1}{\item Vale de ida; y de vuelta solo si se niegan las dos.
  \item Basta: con eso alcanza. Hace falta: sin eso no se puede.
  \item Dos caminos distintos valen si cada paso es cierto.}
\rutade{icf-ma-argumenta-2}{{Uno basta},{Cumplir y fallar},{Siempre o a veces},{Ahora sí, un problema}}
\andamiode{icf-ma-argumenta-2}{\item Mil a favor no la prueban; uno en contra la tumba.
  \item Cumple la primera parte y falla la segunda.
  \item Razón general: siempre. Un caso que falla: a veces.}
\rutade{icf-ma-argumenta-3}{{¿Una sola respuesta?},{Lo que sobra},{Solos o juntos},{Ahora sí, un problema}}
\andamiode{icf-ma-argumenta-3}{\item Dos respuestas posibles: no alcanza.
  \item Cambia el dato: si la respuesta no cambia, sobra.
  \item Uno, el otro y los dos: en ese orden.}
\rutade{icf-ma-finanzas-1}{{Simple: suma lo mismo},{Compuesto: multiplica},{El mismo periodo},{Ahora sí, un problema}}
\andamiode{icf-ma-finanzas-1}{\item Capital por tasa por tiempo: siempre lo mismo.
  \item El interés también gana interés.
  \item Cuenta cuántos periodos de la tasa caben en el tiempo.}
\rutade{icf-ma-finanzas-2}{{El IVA},{Descuento sobre descuento},{La cuota},{Ahora sí, un problema}}
\andamiode{icf-ma-finanzas-2}{\item Ir: por 1.19. Volver: entre 1.19.
  \item Se multiplican las partes que quedan.
  \item Interés de lo que se debe, más lo que se abona.}
\rutade{icf-ma-curvas-1}{{Simetría},{Los ceros},{Máximo o mínimo},{Ahora sí, un problema}}
\andamiode{icf-ma-curvas-1}{\item El eje está entre dos puntos de igual altura.
  \item Producto igual a 0: alguno de los dos es 0.
  \item Ceros, mitad y reemplazo.}
\rutade{icf-ma-curvas-2}{{Por lo mismo},{También baja},{El logaritmo},{Ahora sí, un problema}}
\andamiode{icf-ma-curvas-2}{\item Inicio por base a la x.
  \item Perder 20 \% es quedar con 0.8.
  \item El logaritmo es el exponente que falta.}
\rutade{icf-ma-curvas-3}{{Se repite},{Centro y amplitud},{El seno},{Ahora sí, un problema}}
\andamiode{icf-ma-curvas-3}{\item De cresta a cresta: un periodo.
  \item Centro más y menos amplitud.
  \item El seno va de \ensuremath{-}1 a 1.}
\rutade{icf-ma-curvas-4}{{Una sola salida},{Dominio y rango},{Cortes},{Ahora sí, un problema}}
\andamiode{icf-ma-curvas-4}{\item Una entrada, una salida.
  \item Entra: dominio. Sale: rango.
  \item Eje y: x = 0. Eje x: y = 0.}
\rutade{icf-ma-espacio-1}{{Trasladar y reflejar},{Rotar},{Homotecia},{Ahora sí, un problema}}
\andamiode{icf-ma-espacio-1}{\item Trasladar suma; reflejar cambia un signo.
  \item Gira sin cambiar de tamaño ni de forma.
  \item Lados por k; área por k al cuadrado.}
\rutade{icf-ma-espacio-2}{{Vistas},{Desarrollos},{Otros cuerpos},{Ahora sí, un problema}}
\andamiode{icf-ma-espacio-2}{\item La vista muestra la más alta; el total suma todas.
  \item En fila, salta una.
  \item Cuenta bases y caras.}
\rutade{icf-ma-espacio-3}{{Cartesianas},{Polares},{En el espacio},{Ahora sí, un problema}}
\andamiode{icf-ma-espacio-3}{\item Restar para medir; promediar para el medio.
  \item Giro desde la derecha y pasos.
  \item Derecha, fondo y altura.}
\rutade{icf-ma-dispersion-1}{{Regados o juntos},{Varianza},{Desviación estándar},{Ahora sí, un problema}}
\andamiode{icf-ma-dispersion-1}{\item Rango: el mayor menos el menor.
  \item Restar, elevar, promediar.
  \item Raíz de la varianza.}
\rutade{icf-ma-dispersion-2}{{Percentil},{La caja},{Hallar los cortes},{Ahora sí, un problema}}
\andamiode{icf-ma-dispersion-2}{\item Los de abajo entre el total.
  \item Cuatro tramos, un cuarto en cada uno.
  \item Ordenar, el medio y el medio de cada mitad.}
\rutade{icf-ma-dispersion-3}{{La nube},{Qué tan fuerte},{No es causa},{Ahora sí, un problema}}
\andamiode{icf-ma-dispersion-3}{\item Sube, baja o no tiene rumbo.
  \item Signo: sentido. Cercanía a 1 o \ensuremath{-}1: fuerza.
  \item Busca la tercera cosa.}
\rutade{icf-ma-numeros-1}{{El mismo número},{Sumar y multiplicar},{Notación científica},{Ahora sí, un problema}}
\andamiode{icf-ma-numeros-1}{\item La raya de la fracción es una división.
  \item Sumar: igualar partes. De: multiplicar.
  \item Grandes: positivo. Pequeños: negativo.}
\rutade{icf-ma-numeros-2}{{Cuatro zonas},{Operaciones},{Contar la unión},{Ahora sí, un problema}}
\andamiode{icf-ma-numeros-2}{\item Centro, lunas y afuera.
  \item Alguno, los dos, uno y no el otro, lo demás.
  \item Sumar y quitar una vez lo repetido.}
\rutade{icf-ma-numeros-3}{{Sumar lo mismo},{Multiplicar por lo mismo},{El límite},{Ahora sí, un problema}}
\andamiode{icf-ma-numeros-3}{\item Primero más (lugar \ensuremath{-} 1) saltos.
  \item Primero por razón a la (lugar \ensuremath{-} 1).
  \item El límite es a donde se acerca.}
\rutade{icf-lc-literal-1}{{Volver al texto},{Qué hace cada párrafo},{Dice, se deduce o no está},{Ahora sí, un problema}}
\andamiode{icf-lc-literal-1}{\item La respuesta se señala con el dedo en el texto, no se recuerda.
  \item No preguntes de qué habla un párrafo: pregunta qué hace.
  \item Si no puedes señalar de dónde sale, la opción es inventada.}
\rutade{icf-lc-literal-2}{{La palabra en su lugar},{Hilos hacia atrás},{Bisagras y signos},{Ahora sí, un problema}}
\andamiode{icf-lc-literal-2}{\item Prueba cada opción en la oración: se queda la que no cambia el sentido.
  \item Busca hacia atrás lo que concuerda en género, número y sentido.
  \item El conector dice si la idea choca, resulta o suma; el signo, cómo leerla.}
\rutade{icf-lc-estructura-1}{{El oficio de un párrafo},{Las primeras palabras},{Sostener o adornar},{Ahora sí, un problema}}
\andamiode{icf-lc-estructura-1}{\item Dato, problema, propuesta, objeción, respuesta: cada uno trabaja para el otro.
  \item Lee los comienzos: ahí está el mapa del texto.
  \item Si al quitarla la idea no pierde fuerza, era un adorno.}
\rutade{icf-lc-estructura-2}{{Oponer y limitar},{Causa y consecuencia},{Sumar, ejemplificar, recoger},{Ahora sí, un problema}}
\andamiode{icf-lc-estructura-2}{\item Sin embargo opone, ahora bien limita y aunque reconoce un obstáculo.
  \item Ya que y porque: causa. Por lo tanto y por eso: resultado.
  \item Además suma, por ejemplo muestra, es decir aclara, en resumen recoge.}
\rutade{icf-lc-estructura-3}{{¿Quién habla?},{Directa o indirecta},{Voces que no coinciden},{Ahora sí, un problema}}
\andamiode{icf-lc-estructura-3}{\item Comillas, «según» y «pienso» dicen de quién es cada idea.
  \item Comillas: exacta. «Que»: contada. Y el verbo puede juzgar.
  \item Cada voz escoge sus palabras según lo que quiere mostrar.}
\rutade{icf-lc-argumento-1}{{Tema y tesis},{Razones y adornos},{El otro lado},{Ahora sí, un problema}}
\andamiode{icf-lc-argumento-1}{\item Tesis: una afirmación que se defiende y se puede discutir.
  \item Prueba: tesis + porque + razón. Si no une, no es una razón.
  \item «Se dirá» objeta, «es cierto» concede y «pero» responde.}
\rutade{icf-lc-argumento-2}{{Más firme o más débil},{De pocos a todos},{Dato, ejemplo o imagen},{Ahora sí, un problema}}
\andamiode{icf-lc-argumento-2}{\item Otra causa posible es una de las maneras más fuertes de debilitar.
  \item La conclusión no puede ser más grande que la prueba.
  \item Dato, ejemplo, imagen y autoridad no pesan lo mismo.}
\rutade{icf-lc-literario-1}{{¿Quién cuenta?},{Por lo que hacen},{Lo que cambia},{Ahora sí, un problema}}
\andamiode{icf-lc-literario-1}{\item Yo: dentro. Él o ella: desde afuera. El autor inventa esa voz.
  \item Mira lo que hace y di qué muestra: eso se deduce.
  \item Comienzo, conflicto, decisión y final: el tema está en el cambio.}
\rutade{icf-lc-literario-2}{{Como y es},{Dar vida y exagerar},{Objetos que cargan algo},{Ahora sí, un problema}}
\andamiode{icf-lc-literario-2}{\item Símil: con «como». Metáfora: una cosa es la otra.
  \item Dar vida, exagerar o decir lo contrario: traduce a lo que pasó.
  \item Símbolo: un objeto que vuelve y cambia con los personajes.}
\rutade{icf-lc-literario-3}{{La voz del poema},{Repetir y sonar},{Tema y tono},{Ahora sí, un problema}}
\andamiode{icf-lc-literario-3}{\item La voz poética habla; el poeta la inventa.
  \item La repetición pesa; donde se rompe, está el golpe.
  \item Superficie, fondo y sentimiento: tema y tono.}
\rutade{icf-lc-filosofico-1}{{Una pregunta distinta},{Probar y objetar},{Imaginar para pensar},{Ahora sí, un problema}}
\andamiode{icf-lc-filosofico-1}{\item Si un dato la cierra, no es filosófica.
  \item Pregunta, intento, objeción y nueva respuesta.
  \item «Supongamos que…»: un caso para poner a prueba una idea.}
\rutade{icf-lc-filosofico-2}{{Definir y distinguir},{Ejemplo y analogía},{Objetar y precisar},{Ahora sí, un problema}}
\andamiode{icf-lc-filosofico-2}{\item Busca «llamaré», «entiendo por» y qué separa una idea de otra.
  \item Ejemplo: un caso. Analogía: otra cosa que se le parece.
  \item Objeción, reconocimiento y una idea más precisa.}
\rutade{icf-lc-discontinuo-1}{{Cada parte hace algo},{Tamaño y color},{Números con nota},{Ahora sí, un problema}}
\andamiode{icf-lc-discontinuo-1}{\item Título, tabla, flecha y nota: cada una responde algo distinto.
  \item Grande y rojo atraen; lo pequeño también dice.
  \item Compara con el total: no solo cuenta cuántos, sino de cuántos.}
\rutade{icf-lc-discontinuo-2}{{Globos y cartelas},{Planos y color},{Entre viñetas},{Ahora sí, un problema}}
\andamiode{icf-lc-discontinuo-2}{\item Liso habla, nube piensa, puntas grita, cartela narra.
  \item Lejos empequeñece, cerca emociona, el color marca un cambio.
  \item Une una viñeta con la siguiente: lo que falta se deduce.}
\rutade{icf-lc-discontinuo-3}{{Varias piezas, un tema},{Completan o contradicen},{Concluir con todo},{Ahora sí, un problema}}
\andamiode{icf-lc-discontinuo-3}{\item Pregunta qué se sabría solo con cada pieza.
  \item Una pieza puede mostrar lo que otra deja por fuera.
  \item Si una pieza la contradice, la conclusión dice de más.}
\rutade{icf-lc-evalua-1}{{Lo que se da por hecho},{Para qué escribe},{Recursos y honestidad},{Ahora sí, un problema}}
\andamiode{icf-lc-evalua-1}{\item Supuesto: si lo niegas, la idea se cae.
  \item Pregunta qué haría un lector convencido.
  \item Caso, ironía, duda: cada recurso tiene su efecto.}
\rutade{icf-lc-evalua-2}{{¿Alcanza para decir eso?},{Trampas con nombre},{¿Quién habla?},{Ahora sí, un problema}}
\andamiode{icf-lc-evalua-2}{\item Cuántos, quiénes y comparados con qué.
  \item Costumbre, ataque, falsa elección y generalización.
  \item Interés, límites y letra pequeña.}
\rutade{icf-lc-evalua-3}{{Valores},{Imágenes que delatan},{Acuerdos y choques},{Ahora sí, un problema}}
\andamiode{icf-lc-evalua-3}{\item Mira qué pone primero: ese es su valor.
  \item La imagen destaca algo y esconde lo demás.
  \item Separa el hecho, el valor y la propuesta.}
\rutade{icf-lc-comparar-1}{{Lo que comparten},{Propósito y tono},{Cómo se relacionan},{Ahora sí, un problema}}
\andamiode{icf-lc-comparar-1}{\item «Los dos» exige encontrarlo en cada uno.
  \item Instruir, informar, divertir, convencer, conmover.
  \item Ejemplifica, amplía, matiza o contradice.}
\rutade{icf-lc-comparar-2}{{La tesis de cada uno},{Cómo responde},{Lo que nadie contestó},{Ahora sí, un problema}}
\andamiode{icf-lc-comparar-2}{\item Tesis contra tesis, no tema contra tema.
  \item Conceder, refutar, atacar o usar su propio dato.
  \item Pon cada razón con su respuesta: mira cuál quedó sola.}
\rutade{icf-sc-constitucion-1}{{Estado social de derecho},{Tres ramas},{Los que vigilan},{Ahora sí, un problema}}
\andamiode{icf-sc-constitucion-1}{\item De derecho: todos bajo la ley. Social: el Estado garantiza vida digna.
  \item Legislativa hace leyes, ejecutiva gobierna, judicial decide conflictos.
  \item Contraloría: dinero. Procuraduría: conducta. Defensoría: derechos.}
\rutade{icf-sc-constitucion-2}{{¿De quién es el derecho?},{Las herramientas},{La tutela y su condición},{Ahora sí, un problema}}
\andamiode{icf-sc-constitucion-2}{\item Fundamental: de cada persona. Colectivo: de todos a la vez.
  \item Respuesta, derecho fundamental, derecho de todos, norma o libertad.
  \item Tutela: derecho fundamental en juego y ningún otro medio eficaz.}
\rutade{icf-sc-participacion-1}{{Elegir y decidir},{Cada mecanismo, su uso},{Retirar del cargo},{Ahora sí, un problema}}
\andamiode{icf-sc-participacion-1}{\item Representativa: eligen y otros deciden. Participativa: deciden los ciudadanos.
  \item Norma: referendo. Presidente: plebiscito. Asunto: consulta. Hablar: cabildo.
  \item Solo alcaldes y gobernadores, y después de un año.}
\rutade{icf-sc-participacion-2}{{Dinero, conducta, delito},{¿Dónde están?},{Los ojos de la gente},{Ahora sí, un problema}}
\andamiode{icf-sc-participacion-2}{\item Contraloría: dinero. Procuraduría: conducta. Fiscalía: delito.
  \item La Fiscalía investiga delitos, pero es de la rama judicial.
  \item La veeduría pide, mira y denuncia; sancionar es del Estado.}
\rutade{icf-sc-participacion-3}{{La regla mayor},{Tres caminos},{¿Por qué tan difícil?},{Ahora sí, un problema}}
\andamiode{icf-sc-participacion-3}{\item Ninguna norma puede contradecir a la Constitución.
  \item Acto legislativo, referendo o asamblea constituyente.
  \item La dificultad protege a quien hoy es minoría.}
\rutade{icf-sc-perspectivas-1}{{Posición e interés},{Muchos actores},{Mismo lado, distinta razón},{Ahora sí, un problema}}
\andamiode{icf-sc-perspectivas-1}{\item Posición: qué pide. Interés: para qué lo pide.
  \item Pregunta quién gana, quién pierde, quién decide y quién falta.
  \item Coincidir no es pensar igual, y un grupo no es una sola voz.}
\rutade{icf-sc-perspectivas-2}{{El lente},{Prejuicio},{La misma acción},{Ahora sí, un problema}}
\andamiode{icf-sc-perspectivas-2}{\item Cosmovisión: la manera de ver el mundo desde la que se actúa.
  \item Todo un grupo, sin conocerlo y sin datos: prejuicio.
  \item Busca la razón que la persona da, no la que tú supondrías.}
\rutade{icf-sc-perspectivas-3}{{¿A quién deja por fuera?},{¿Se puede hacer?},{¿La van a cumplir?},{Ahora sí, un problema}}
\andamiode{icf-sc-perspectivas-3}{\item No hay solución sin alguien que pierda algo: búscalo.
  \item Viable: alcanza el dinero, el tiempo y hay quién la cumpla.
  \item Acordada, justa y que se pueda comprobar.}
\rutade{icf-sc-fuentes-1}{{Primaria o secundaria},{Quién, para quién, para qué},{El contexto y el silencio},{Ahora sí, un problema}}
\andamiode{icf-sc-fuentes-1}{\item Primaria: de la época. Secundaria: la estudia. Ninguna es verdad automática.
  \item El propósito explica lo que una fuente muestra y lo que calla.
  \item Contexto para entender; silencio para saber qué más buscar.}
\rutade{icf-sc-fuentes-2}{{Qué mide cada una},{Los límites},{Comparar y medir las palabras},{Ahora sí, un problema}}
\andamiode{icf-sc-fuentes-2}{\item Pregunta qué se contó exactamente: solo eso queda probado.
  \item Busca a quién deja por fuera la fuente.
  \item Compara, di lo que se sabe y busca otra explicación.}
\rutade{icf-sc-fuentes-3}{{Leer una caricatura},{La ironía},{Leer un discurso},{Ahora sí, un problema}}
\andamiode{icf-sc-fuentes-3}{\item Qué agranda, qué simboliza y qué contrasta: ahí está la opinión.
  \item Ironía: la frase choca con la escena.
  \item Quita ataques y falsas elecciones: lo que queda se revisa.}
\rutade{icf-sc-sistemico-1}{{Detonante y fondo},{En cadena y en círculo},{Juntas no es por},{Ahora sí, un problema}}
\andamiode{icf-sc-sistemico-1}{\item El detonante dispara; las causas de fondo explican.
  \item Sigue las flechas: algunas vuelven al comienzo.
  \item Antes de culpar, busca qué más cambió y compara casos.}
\rutade{icf-sc-sistemico-2}{{Las dimensiones},{Chocan},{Ceder en cada una},{Ahora sí, un problema}}
\andamiode{icf-sc-sistemico-2}{\item Económica, ambiental, cultural, legal, política, social.
  \item Di qué mejora y qué empeora: ese es el choque.
  \item La salida que atiende varias dimensiones dura más.}
\rutade{icf-sc-sistemico-3}{{Lo que no se busca},{Para quién},{Medir y corregir},{Ahora sí, un problema}}
\andamiode{icf-sc-sistemico-3}{\item Efecto no buscado: lo que produce la gente al acomodarse.
  \item Pregunta a quién le pesa más y qué esconde el promedio.
  \item Mide la meta y decide antes cuándo corregir.}
\rutade{icf-sc-historia-1}{{El orden colonial},{Resistir},{Cambios y permanencias},{Ahora sí, un problema}}
\andamiode{icf-sc-historia-1}{\item Cargos, comercio y vida según el origen.
  \item Palenques, pleitos, marchas y memoria.
  \item Una fecha no cambia todo: mira qué siguió.}
\rutade{icf-sc-historia-2}{{Tres preguntas},{Dos partidos},{El péndulo},{Ahora sí, un problema}}
\andamiode{icf-sc-historia-2}{\item Centro o regiones, Iglesia y voto.
  \item En general: centro e Iglesia; regiones y libertades.
  \item Rionegro federal, 1886 central, y el café al fondo.}
\rutade{icf-sc-historia-3}{{De la huelga al 9 de abril},{Un pacto con puerta cerrada},{Víctimas y paz},{Ahora sí, un problema}}
\andamiode{icf-sc-historia-3}{\item La violencia venía de antes; el 9 de abril la multiplicó.
  \item Menos guerra entre dos; menos espacio para los demás.
  \item Reconocer a las víctimas y contar también los intentos de paz.}
\rutade{icf-sc-mundo-1}{{Dos clases de revolución},{¿Derechos para todos?},{Fábricas y nuevas ideas},{Ahora sí, un problema}}
\andamiode{icf-sc-mundo-1}{\item Políticas: quién manda. Industrial: cómo se produce.
  \item Compara lo que se proclama con quién queda por fuera.
  \item Del conflicto salen sindicatos, leyes e ideas.}
\rutade{icf-sc-mundo-2}{{La chispa y la pólvora},{De una guerra a otra},{Una guerra fría… para ellas},{Ahora sí, un problema}}
\andamiode{icf-sc-mundo-2}{\item Colonias, alianzas y nacionalismo; Sarajevo, la chispa.
  \item Humillación, crisis, dictadura, guerra y genocidio.
  \item Fría entre potencias, caliente en otros países.}
\rutade{icf-sc-mundo-3}{{Dejar de ser colonia},{Rayas y leyes},{Movimientos sociales},{Ahora sí, un problema}}
\andamiode{icf-sc-mundo-3}{\item Dos causas y dos caminos.
  \item La independencia no borra las rayas ni las leyes.
  \item Organizarse, fuera de las instituciones, para cambiar algo.}
\rutade{icf-sc-geografia-1}{{Regiones y altura},{Mapas que escogen},{La población se mueve},{Ahora sí, un problema}}
\andamiode{icf-sc-geografia-1}{\item Seis regiones; al subir mil metros, unos seis grados menos.
  \item Qué muestra, qué deja fuera y cuánto redujo.
  \item Densidad, migración y pirámide.}
\rutade{icf-sc-geografia-2}{{Renovable no es infinito},{Tres patas},{Quién está en la mesa},{Ahora sí, un problema}}
\andamiode{icf-sc-geografia-2}{\item Se recupera si se usa despacio.
  \item Si falla una, no es sostenible.
  \item Pregunta quién gana, quién pierde y quién decide.}
\rutade{icf-sc-economia-1}{{Oferta y demanda},{Quién manda},{Impuestos},{Ahora sí, un problema}}
\andamiode{icf-sc-economia-1}{\item Menos oferta o más demanda: sube el precio.
  \item Competencia, monopolio y cartel.
  \item Directo o indirecto; progresivo o regresivo.}
\rutade{icf-sc-economia-2}{{Inflación},{Quién es desempleado},{Pobreza y reparto},{Ahora sí, un problema}}
\andamiode{icf-sc-economia-2}{\item Nominal es en pesos; real, en compras.
  \item Buscar y no encontrar.
  \item Lo que falta y cómo se reparte.}
\rutade{icf-sc-economia-3}{{Una red},{Abrir la puerta},{También viajan},{Ahora sí, un problema}}
\andamiode{icf-sc-economia-3}{\item Comercio, dinero, información y personas.
  \item La puerta se abre en los dos sentidos.
  \item Lo bueno y lo malo cruzan fronteras.}
\rutade{icf-sc-modelos-1}{{Tres que se confunden},{Poder y legitimidad},{Mirar a la gente},{Ahora sí, un problema}}
\andamiode{icf-sc-modelos-1}{\item Organización, personas y sentimiento de grupo.
  \item Legal: según la norma. Legítimo: reconocido.
  \item Capas que se cruzan y culturas que se entienden.}
\rutade{icf-sc-modelos-2}{{Un modelo},{Cuatro lógicas},{Evaluar una decisión},{Ahora sí, un problema}}
\andamiode{icf-sc-modelos-2}{\item Cuatro preguntas para comparar.
  \item Libertad, igualdad, propiedad común y comunidad.
  \item Qué valor, quién gana y quién queda por fuera.}
\rutade{icf-cn-indaga-1}{{Cambiar una sola cosa},{Los nombres de las variables},{Una hipótesis se puede tumbar},{Ahora sí, un problema}}
\andamiode{icf-cn-indaga-1}{\item Si cambian dos cosas a la vez, no se sabe cuál fue.
  \item Independiente: se cambia. Dependiente: se mide. Controlada: igual.
  \item Una hipótesis dice algo que se puede medir y que puede resultar falso.}
\rutade{icf-cn-indaga-2}{{Leer la tendencia},{Qué se puede concluir},{Repetir y medir bien},{Ahora sí, un problema}}
\andamiode{icf-cn-indaga-2}{\item Mira si sube, baja o sigue igual, y cuánto en cada paso.
  \item Si habla de algo que no se midió, dice de más.
  \item Una sola medida no muestra su error; varias, sí.}
\rutade{icf-cn-celula-1}{{El agua busca lo salado},{A favor o en contra},{Con pared o sin pared},{Ahora sí, un problema}}
\andamiode{icf-cn-celula-1}{\item Ósmosis: el agua va hacia donde hay más sustancia disuelta.
  \item A favor es gratis; en contra hace falta energía.
  \item Con pared, la célula se pone rígida; sin ella, puede reventar.}
\rutade{icf-cn-celula-2}{{Volver al equilibrio},{El cuerpo se avisa},{Nadie trabaja solo},{Ahora sí, un problema}}
\andamiode{icf-cn-celula-2}{\item Homeostasis: ante un cambio, el cuerpo vuelve a su valor.
  \item Cambio, aviso, respuesta y efecto que lo corrige.
  \item Romper, repartir, respirar y limpiar: un solo equipo.}
\rutade{icf-cn-celula-3}{{Copiar la receta},{Leer de tres en tres},{Cuando cambia una letra},{Ahora sí, un problema}}
\andamiode{icf-cn-celula-3}{\item Transcripción: A con U, T con A, G con C y C con G.
  \item Tres letras, un codón; un codón, un aminoácido.
  \item Cambiar una letra afecta un codón; quitar una, todos los que siguen.}
\rutade{icf-cn-movimiento-1}{{La gráfica cuenta un viaje},{Más rápido o más lento},{Media y signo},{Ahora sí, un problema}}
\andamiode{icf-cn-movimiento-1}{\item Recta: velocidad constante. Horizontal: quieto. Baja: se devuelve.
  \item Velocidad: cuánto cambia la posición. Aceleración: cuánto cambia la velocidad.
  \item Media: total entre total. Signo: hacia dónde.}
\rutade{icf-cn-movimiento-2}{{Juntar las fuerzas},{Fuerza, masa y aceleración},{Fuerzas por pares},{Ahora sí, un problema}}
\andamiode{icf-cn-movimiento-2}{\item Mismo sentido se suma; contrario se resta. Neta cero: no cambia.
  \item Más fuerza, más aceleración; más masa, menos.
  \item Acción y reacción: iguales, contrarias y sobre cuerpos distintos.}
\rutade{icf-cn-movimiento-3}{{Caer igual},{Masa y peso},{Orbitar es caer},{Ahora sí, un problema}}
\andamiode{icf-cn-movimiento-3}{\item Sin aire, la masa no cambia la caída: unos 10 m/s más cada segundo.
  \item La masa no cambia de lugar a lugar; el peso, sí.
  \item La Luna sí cae: cae alrededor de la Tierra.}
\rutade{icf-cn-materia-1}{{Partículas},{El tramo plano},{Densidad},{Ahora sí, un problema}}
\andamiode{icf-cn-materia-1}{\item Cambia el movimiento y la distancia; no las partículas.
  \item Plano: cambio de estado. A menor presión, hierve a menos.
  \item Densidad: masa entre volumen. Lo menos denso flota.}
\rutade{icf-cn-materia-2}{{Mezcla o compuesto},{Busca la diferencia},{El orden y la masa},{Ahora sí, un problema}}
\andamiode{icf-cn-materia-2}{\item Mezcla: juntas. Compuesto: unidas. Homogénea: pareja.
  \item Imán, filtro, decantar, evaporar o destilar: según la diferencia.
  \item Nada se pierde; el porcentaje va sobre el total.}
\rutade{icf-cn-materia-3}{{Hasta cuánto cabe},{Más rápido o más},{Gases y semejantes},{Ahora sí, un problema}}
\andamiode{icf-cn-materia-3}{\item Menos del límite: insaturada. En el límite: saturada. Lo demás, al fondo.
  \item Agitar y moler apuran. Calentar o agregar agua: cabe más.
  \item Más calor, menos gas. Lo semejante disuelve a lo semejante.}
\rutade{icf-cn-ecosistemas-1}{{La cadena},{La décima parte},{Mover un eslabón},{Ahora sí, un problema}}
\andamiode{icf-cn-ecosistemas-1}{\item Productor, primario, secundario, terciario; y descomponedores para todos.
  \item Cada paso, cerca del 10 \%: por eso arriba hay pocos.
  \item Mira quién lo comía y a quién se comía.}
\rutade{icf-cn-ecosistemas-2}{{El carbono},{El nitrógeno},{Demasiado},{Ahora sí, un problema}}
\andamiode{icf-cn-ecosistemas-2}{\item Fotosíntesis toma; respirar, descomponer y quemar devuelven.
  \item Aire, bacterias, suelo y raíz.
  \item Nitratos, algas, descomposición y falta de oxígeno.}
\rutade{icf-cn-ecosistemas-3}{{Quién gana y quién pierde},{Suben y bajan},{El límite},{Ahora sí, un problema}}
\andamiode{icf-cn-ecosistemas-3}{\item Gana-gana, gana-igual, gana-pierde y pierde-pierde.
  \item Primero las presas, después los depredadores.
  \item Capacidad de carga: cuántos sostiene el lugar.}
\rutade{icf-cn-energia-1}{{Dos formas},{Con roce},{El cuadrado},{Ahora sí, un problema}}
\andamiode{icf-cn-energia-1}{\item Sin roce, potencial más cinética siempre da lo mismo.
  \item La energía total se conserva; la mecánica, solo sin roce.
  \item Doble velocidad: cuatro veces. Doble masa: dos veces.}
\rutade{icf-cn-energia-2}{{No son lo mismo},{Tres caminos},{Unos cuestan más},{Ahora sí, un problema}}
\andamiode{icf-cn-energia-2}{\item Temperatura se tiene; calor pasa, hasta el equilibrio.
  \item Contacto, fluido u ondas; el aire quieto frena.
  \item Mayor calor específico: cambia más despacio.}
\rutade{icf-cn-atomo-1}{{Tres partículas},{Isótopos e iones},{La tabla},{Ahora sí, un problema}}
\andamiode{icf-cn-atomo-1}{\item Protones: qué elemento es. Masa: protones más neutrones.
  \item Isótopo: otros neutrones. Ion: otros electrones.
  \item Grupo: se parecen. Bajar: crece. Avanzar: se encoge.}
\rutade{icf-cn-atomo-2}{{Entregar},{Compartir},{Un mar},{Ahora sí, un problema}}
\andamiode{icf-cn-atomo-2}{\item Iónico: metal con no metal, red dura, conduce disuelto.
  \item Covalente: comparten. Igual, apolar; desigual, polar.
  \item Metálico: conduce sólido y se dobla.}
\rutade{icf-cn-herencia-1}{{Dos alelos},{El cuadro},{Lo que se espera},{Ahora sí, un problema}}
\andamiode{icf-cn-herencia-1}{\item Letras: genotipo. Lo que se ve: fenotipo.
  \item Bb con Bb: tres a uno.
  \item Mientras más nacimientos, más se acerca.}
\rutade{icf-cn-herencia-2}{{Uno o dos progenitores},{Dos divisiones},{El seguro de la variedad},{Ahora sí, un problema}}
\andamiode{icf-cn-herencia-2}{\item Asexual: uno, iguales. Sexual: dos, variados.
  \item Mitosis: dos iguales. Meiosis: cuatro con la mitad.
  \item Variedad: el seguro de vida de una especie.}
\rutade{icf-cn-herencia-3}{{La población cambia},{Al azar y escogido},{Las pruebas},{Ahora sí, un problema}}
\andamiode{icf-cn-herencia-3}{\item Variación, herencia, más nacen y unos se reproducen más.
  \item Primero la variación; después, la selección.
  \item Roca, huesos, ADN y lo que pasa hoy.}
\rutade{icf-cn-reacciones-1}{{Físico o químico},{Los átomos se reacomodan},{La masa no cambia},{Ahora sí, un problema}}
\andamiode{icf-cn-reacciones-1}{\item Misma sustancia: físico. Sustancia nueva: químico.
  \item Los mismos átomos antes y después.
  \item Entra lo mismo que sale.}
\rutade{icf-cn-reacciones-2}{{El mol},{La receta},{El que limita},{Ahora sí, un problema}}
\andamiode{icf-cn-reacciones-2}{\item Gramos entre masa molar: moles.
  \item Cuántas veces se hace la receta.
  \item Divide lo que hay entre lo que pide la receta.}
\rutade{icf-cn-reacciones-3}{{La escala},{Por diez},{Neutralizar y diluir},{Ahora sí, un problema}}
\andamiode{icf-cn-reacciones-3}{\item Menos de 7 ácida; más de 7 básica.
  \item Cuenta los pasos: un diez por cada uno.
  \item Con base se puede pasar de 7; con agua, no.}
\rutade{icf-cn-ondas-1}{{Viaja la energía},{Medir una onda},{La rapidez},{Ahora sí, un problema}}
\andamiode{icf-cn-ondas-1}{\item Viaja la energía, no la materia.
  \item Altura, distancia y cuántas por segundo.
  \item Si una sube, la otra baja.}
\rutade{icf-cn-ondas-2}{{Rebotar},{El eco},{Doblarse},{Ahora sí, un problema}}
\andamiode{icf-cn-ondas-2}{\item El ángulo que llega es igual al que sale.
  \item Recorrido entre dos.
  \item Más despacio, se acerca; más rápido, se aleja.}
\rutade{icf-cn-gases-1}{{Boyle},{Charles},{Recipiente rígido},{Ahora sí, un problema}}
\andamiode{icf-cn-gases-1}{\item Menos volumen, más presión.
  \item Volumen y kelvin, en proporción.
  \item Mira qué queda fijo.}
\rutade{icf-cn-gases-2}{{Soluto y disolución},{Diluir},{Otras cosas que se hacen},{Ahora sí, un problema}}
\andamiode{icf-cn-gases-2}{\item Soluto entre disolución.
  \item Gramos iguales; resta lo que ya había.
  \item Pregunta qué salió y qué entró.}
\rutade{icf-cn-electrico-1}{{Ohm},{Serie y paralelo},{Potencia},{Ahora sí, un problema}}
\andamiode{icf-cn-electrico-1}{\item Corriente es voltaje entre resistencia.
  \item Serie: un camino. Paralelo: varios.
  \item Vatios por horas.}
\rutade{icf-cn-electrico-2}{{Cargas},{Imanes},{De ida y vuelta},{Ahora sí, un problema}}
\andamiode{icf-cn-electrico-2}{\item Al doble de lejos, la cuarta parte.
  \item Dos polos, siempre.
  \item Ida: electroimán. Vuelta: inducción.}
\rutade{icf-cn-cts-1}{{Prevenir y curar},{Comparar riesgos},{Ver por dentro},{Ahora sí, un problema}}
\andamiode{icf-cn-cts-1}{\item Prevenir no es curar.
  \item Divide cada caso entre su grupo.
  \item Beneficio frente a riesgo.}
\rutade{icf-cn-cts-2}{{Energía},{Efecto invernadero},{Residuos},{Ahora sí, un problema}}
\andamiode{icf-cn-cts-2}{\item Ventajas y costos de principio a fin.
  \item Necesario, pero sin exceso.
  \item En ese orden.}
\rutade{icf-cn-cts-3}{{Testimonio y prueba},{Un buen experimento},{Datos y valores},{Ahora sí, un problema}}
\andamiode{icf-cn-cts-3}{\item Sincero no es lo mismo que probado.
  \item Cambiar una sola cosa.
  \item Datos primero, y después valores.}
\rutade{icf-en-avisos-1}{{¿Dónde está el aviso?},{¿Ordena o informa?},{Decirlo con otras palabras},{Ahora sí, un problema}}
\andamiode{icf-en-avisos-1}{\item Busca la palabra que solo tiene sentido en un lugar.
  \item Must obliga, no -ing prohíbe, please pide; lo demás informa.
  \item Desconfía de la opción que copia palabras: compara la idea.}
\rutade{icf-en-avisos-2}{{¿Persona, lugar o cosa?},{Las pistas del verbo},{Palabras que apuntan},{Ahora sí, un problema}}
\andamiode{icf-en-avisos-2}{\item Primero la clase de palabra, después el significado.
  \item -ed ya pasó, will pasará, has… for sigue hasta hoy.
  \item He, she: personas. It: una cosa. They, them: varias.}
\rutade{icf-en-dialogos-1}{{La palabra que pregunta},{El auxiliar se repite},{La trampa de la palabra repetida},{Ahora sí, un problema}}
\andamiode{icf-en-dialogos-1}{\item Where lugar, when momento, who persona, how long duración, why razón.
  \item Do con do, is con is, can con can, did con did.
  \item Busca el dato que se pide, no la palabra que se repite.}
\rutade{icf-en-dialogos-2}{{Pedir y ofrecer},{Do you mind?},{Quejas, gracias y disculpas},{Ahora sí, un problema}}
\andamiode{icf-en-dialogos-2}{\item Could you pide, would you like ofrece, may I pide permiso.
  \item Con «Do you mind», decir no es dar permiso.
  \item Queja, gracias, invitación o disculpa: cada una con su fórmula.}
\rutade{icf-en-dialogos-3}{{Cómo se siente},{Planes y promesas},{¿De acuerdo?},{Ahora sí, un problema}}
\andamiode{icf-en-dialogos-3}{\item Triste se lamenta, un logro se felicita, un susto se anima.
  \item Going to: plan. Will: promesa o cálculo. What if: sugerencia.
  \item So do I y neither do I coinciden; I’d rather prefiere otra cosa.}
\rutade{icf-en-gramatica-1}{{¿Terminó o llega hasta hoy?},{Lo que estaba pasando},{Mañana},{Ahora sí, un problema}}
\andamiode{icf-en-gramatica-1}{\item Yesterday, ago, in 2015: pasado. Since y for: have + participio.
  \item Larga: was o were + -ing. Corta: pasado simple.
  \item Cita: am -ing. Plan: going to. Predicción: will. If + presente.}
\rutade{icf-en-gramatica-2}{{In, on, at},{¿Se puede contar?},{Unir ideas},{Ahora sí, un problema}}
\andamiode{icf-en-gramatica-2}{\item In: mes o adentro. On: día o encima. At: hora o punto.
  \item Cuenta: many, a few. No se cuenta: much, a little.
  \item So: resultado. Because: causa. Although y but: contraste.}
\rutade{icf-en-gramatica-3}{{Más que},{Igual o menos},{El más de todos},{Ahora sí, un problema}}
\andamiode{icf-en-gramatica-3}{\item Corta: -er than. Larga: more… than. Nunca las dos juntas.
  \item As… as: igual. Not as… as: menos.
  \item The + -est o the most. Good: best. Bad: worst.}
\rutade{icf-en-corto-1}{{¿Quién le escribe a quién?},{¿Para qué?},{Los datos y lo que se entiende},{Ahora sí, un problema}}
\andamiode{icf-en-corto-1}{\item Saludo: quien recibe. Firma: quien escribe.
  \item Pide, avisa, invita o agradece.
  \item Qué, dónde, cuándo; y une dos datos para entender más.}
\rutade{icf-en-corto-2}{{La idea de todo el texto},{Correct, wrong, doesn't say},{La palabra por el contexto},{Ahora sí, un problema}}
\andamiode{icf-en-corto-2}{\item La idea principal cubre todo; un detalle, solo una parte.
  \item Correct: está. Wrong: dice lo contrario. Doesn't say: no está.
  \item Lee antes y después: el contexto la explica.}
\rutade{icf-en-largo-1}{{Primero, de qué trata},{La misma idea, otras palabras},{Detalle o idea},{Ahora sí, un problema}}
\andamiode{icf-en-largo-1}{\item Título y primeras oraciones: un mapa del texto.
  \item Used to: antes. No longer y not any more: ya no.
  \item La principal cubre todo; el detalle, una parte.}
\rutade{icf-en-largo-2}{{¿Para qué escribe?},{Hecho u opinión},{Aceptar y limitar},{Ahora sí, un problema}}
\andamiode{icf-en-largo-2}{\item To advise, to complain, to explain, to tell.
  \item Fact: se comprueba. Opinion: alguien puede no estar de acuerdo.
  \item «I understand… However… I am asking for…».}
\rutade{icf-en-vocabulario-1}{{Mira alrededor},{Cada palabra y su compañía},{Verbos con partícula},{Ahora sí, un problema}}
\andamiode{icf-en-vocabulario-1}{\item La palabra debe encajar en sentido y en gramática.
  \item Make crea, do realiza, say algo, tell a alguien.
  \item La partícula manda: after cuida, for busca, down rechaza.}
\rutade{icf-en-vocabulario-2}{{Familias},{Prefijos y -ed o -ing},{Conectores},{Ahora sí, un problema}}
\andamiode{icf-en-vocabulario-2}{\item -ful y -ous: cómo es. -ly: cómo se hace. -tion: la cosa.
  \item Un- y dis-: contrario. -ed siente, -ing causa.
  \item However contrasta, moreover suma, therefore resulta.}
\rutade{icf-en-inferencia-1}{{Unir dos datos},{Los gestos},{Must, might, can't},{Ahora sí, un problema}}
\andamiode{icf-en-inferencia-1}{\item Dos datos señalados, una idea.
  \item Gesto, razón, emoción.
  \item Must: casi seguro. Might: quizá. Can't: imposible.}
\rutade{icf-en-inferencia-2}{{Fact or opinion},{Opinion + reason},{Agree, disagree, change},{Ahora sí, un problema}}
\andamiode{icf-en-inferencia-2}{\item Contar o comprobar: fact. Should, I think: opinion.
  \item Opinion + because + reason.
  \item You are right… but… changed my mind.}
\rutade{pu-ap-1}{{Sacar el principio},{La estructura, no el tema},{Del caso al límite},{Ahora sí, un problema}}
\andamiode{pu-ap-1}{\item El principio dice lo mismo sin nombrar el caso, y no exagera.
  \item Busca las piezas del principio en el caso. Si falta una, no se cumple.
  \item «Va más allá» saca la idea; «eso no significa» le pone un límite.}
\rutade{pu-ap-2}{{La misma relación},{Tipos de relación},{Hasta dónde llega},{Ahora sí, un problema}}
\andamiode{pu-ap-2}{\item Di la relación con un verbo y prueba el mismo verbo al otro lado.
  \item Parte y todo, herramienta y usuario, causa y efecto, etapa y lo que llega a ser.
  \item «Así como… también» dice en qué se parecen; «pero», en qué no.}
\rutade{pu-fu-1}{{Directa e indirecta},{Quién habla y quién escribe},{Lo que concluye el autor},{Ahora sí, un problema}}
\andamiode{pu-fu-1}{\item Comillas: palabras exactas. «Que» o «según»: contadas por quien escribe.
  \item «Yo», «creo», «me parece», «a mi modo de ver» son de quien escribe. Lo demás tiene otro dueño.
  \item «No… sino…» suele decir la tesis. Un dato la apoya, la debilita o no pesa.}
\rutade{pu-fu-2}{{Tres posturas},{La cita de afuera},{Cómo se matiza},{Ahora sí, un problema}}
\andamiode{pu-fu-2}{\item A favor, en contra o sí, pero con condiciones.
  \item Supón que la cita es cierta: ¿la tesis queda más fuerte, más débil o igual?
  \item «Aunque», «pero», «reconozco»: matiz. «Propone» no es «hizo» ni «la mayor parte» es «todo».}
\rutade{pu-cn-1}{{Lo que anuncia un conector},{Oponer y rectificar},{Lo que necesita un conector},{Ahora sí, un problema}}
\andamiode{pu-cn-1}{\item Contraste, consecuencia, suma, ejemplo, causa y conclusión: el conector dice qué esperar.
  \item Oponer deja dos ideas en pie; rectificar niega una y pone otra en su lugar.
  \item Cada conector pide algo antes. Después del contraste viene lo que más pesa.}
\rutade{pu-cn-2}{{Lo que hace un párrafo},{Matizar y relativizar},{Los signos},{Ahora sí, un problema}}
\andamiode{pu-cn-2}{\item Ejemplificar, conceder, refutar, concluir: busca la marca en el comienzo.
  \item ¿De quién es la idea que se limita? Propia: matiza. Ajena: relativiza.
  \item Comillas irónicas, rayas que meten ejemplos, paréntesis que guardan condiciones.}
\rutade{pu-ag-1}{{Qué y por qué},{Tipos de argumento},{La objeción y su respuesta},{Ahora sí, un problema}}
\andamiode{pu-ag-1}{\item Tesis: ¿qué defiende? Argumento: ¿por qué?
  \item Datos, ejemplo, autoridad, analogía y causa.
  \item «Hay quien piensa» objeta, «sin duda» concede y un hecho puede responder sin «pero».}
\rutade{pu-ag-2}{{Para qué escribe},{El tono},{Comillas con intención},{Ahora sí, un problema}}
\andamiode{pu-ag-2}{\item Informar, persuadir, criticar o reflexionar: ¿qué quiere lograr en quien lee?
  \item Quita el adverbio: si el dato sigue, lo que perdiste era la actitud del autor.
  \item Comillas que toman distancia o muestran lo que falta. Y criticar algo no es burlarse de las personas.}
\rutade{pu-ag-3}{{Comprobar o valorar},{Opiniones escondidas},{Hechos que apoyan opiniones},{Ahora sí, un problema}}
\andamiode{pu-ag-3}{\item Pregunta: si alguien lo duda, ¿hay a dónde ir a mirar?
  \item No basta la palabra: «rico en» describe y «ya es hora» opina.
  \item Para probar hace falta un hecho; la opinión se apoya en él.}
\rutade{pu-dt-1}{{En los dos o en uno},{Uno explica al otro},{Con qué se sostiene},{Ahora sí, un problema}}
\andamiode{pu-dt-1}{\item Busca cada idea en los dos. Si solo está en uno, es diferencia.
  \item Pon «porque» entre los dos: si tiene sentido, uno explica al otro.
  \item Investigación, lo que cuenta otro o experiencia propia: mira la persona y la cautela.}
\rutade{pu-dt-2}{{La tesis como lente},{Conceder, proponer, rechazar},{Un tercero en la discusión},{Ahora sí, un problema}}
\andamiode{pu-dt-2}{\item Pon la frase frente a la tesis: si la contradice, no lo diría.
  \item Lo que concede cuenta como suyo; lo que «los críticos dicen» y él rechaza, no.
  \item Sirve lo que su autor sostiene o concede y choca con la tesis.}
\rutade{pu-rf-1}{{El referente},{Este y aquel},{Lo que da por contado},{Ahora sí, un problema}}
\andamiode{pu-rf-1}{\item Busca hacia atrás, reemplaza y lee: si tiene sentido, ese es.
  \item Este, el cercano; aquel, el lejano; esto último, el final de la lista.
  \item «Un» presenta, «el» retoma; «en cambio» y «al día siguiente» van después de algo.}
\rutade{pu-rf-2}{{Otro nombre para lo mismo},{Lo que se calla},{Contar en orden y unir},{Ahora sí, un problema}}
\andamiode{pu-rf-2}{\item Sinónimo, palabra más general o apodo: nombres distintos para lo mismo.
  \item Lo que falta está en la primera mitad; la coma guarda su puesto.
  \item Cuenta en orden y confirma con el género y el número.}
\rutade{pu-lf-1}{{Cuatro figuras},{Qué quiere decir},{Dónde se esconde},{Ahora sí, un problema}}
\andamiode{pu-lf-1}{\item Es: metáfora. Como: comparación. Humano: personificación. Exagera: hipérbole.
  \item Pregunta qué hay de cierto detrás de la imagen y dilo sin agregar.
  \item La coma puede abrir una metáfora; los dos puntos, su razón. En el resumen no van.}
\rutade{pu-lf-2}{{Decir lo contrario},{La ironía sutil},{Lo que de verdad piensa},{Ahora sí, un problema}}
\andamiode{pu-lf-2}{\item Si lo que se dice choca con lo que pasa, es ironía.
  \item La sutil no dice lo contrario de un hecho: pone un tono que no va con el asunto.
  \item Dale la vuelta al elogio: ahí está la tesis. Los datos van en serio.}
